Performance varied substantially across mechanisms. The G1 near-budget F1 advantage over Variance did not transfer: in G2 the point estimate favors Variance, while the SU40-minus-Variance F1 interval spans zero, and its Recall interval is [-14.58, -2.97] percentage points. Fixed20+C11 also has higher point F1 (55.10\%) at an actual sample gap of -2.65\%; no predeclared fixed-point interval was estimated. Thus the data do not support SU40 quality superiority at equal acquisition cost. All fixed points remain visible, including Fixed10's F1 69.98\%. SU40 instead uses -67.92\% bytes and -56.51\% operational fault-induced reports relative to Variance, with a larger ordinary-data age. The U diagnostic on G1 and Gate+U on G2 show that part of this advantage is reusable communication behavior, not unique sampling novelty. SU20 improves G2 Recall over SU40, but sample cycles change by +43.84\% and FI counts are 193 versus 157. Its short-event Recall remains only 37.34\%, despite 92.34\% sampled coverage. SU40 short Recall (9.22\%) is below GateOnly (9.69\%), another reversal of the G1 short-event pattern. Sampling coverage alone does not guarantee debounce, gating or shape-specific detection success. Rao-inspired remains competitive on quality despite fixed-center mismatch, but its actual sample budget grows well beyond SU40; this is neither matched-budget evidence nor a faithful Rao reproduction. The contrast also combines changed signal shapes, correlated noise, amplitudes, timing and label geometry, so these data cannot uniquely attribute a reversal to one causal factor. The revised contribution is a reproducible conditional systems tradeoff and falsification of an overly broad transfer claim, not evidence of a universally preferable balanced mode.
